% =====================================================================
% Section 6 — Results and Comparisons
%
% A consolidated results section for the two microcode case studies
% (Keccak-f[1600] and Curve25519 field arithmetic), written in the style
% and to the comprehensiveness of the "Results and Comparisons" section of
% Zhang, Yan, Huang and Koc, "Optimized Software Implementation of Keccak,
% Kyber, and Dilithium on RV{32,64}IM{B}{V}", IACR TCHES 2025(1):632-655
% (referred to below as Zhang et al.\ [TCHES 2025(1)]).
%
% All numbers are taken verbatim from the committed measurement files:
%   curve25519/RESULTS.md          (X25519 end-to-end + ratio matrices)
%   keccak/bench/KECCAK_RESULTS.md (Keccak head-to-head matrix, 13 contenders)
% and from the per-op microbenchmarks reported in Sections 4-5. The stability
% figures are read off the per-configuration min matrix in KECCAK_RESULTS.md.
%
% PLATFORM FACTS (verified 2026-06-28, see section_challenges.tex):
%   Goldmont (Celeron N3350) is OUT-OF-ORDER, 3-wide, 3 integer ALU ports.
%   Its ISA reaches SSE4.2/AES/SHA but has NO BMI2 (MULX), NO ADX
%   (ADCX/ADOX), and NO AVX/AVX2. Do not describe it as in-order.
%
% Cross-references (\Cref{sec:curve25519}, \Cref{sec:keccak-design},
% \Cref{sec:findings}, \Cref{sec:budget}, \Cref{sec:ucode-layer}) point at
% the case-study and microcode-layer sections. Requires booktabs.
% NOTE: this file and section_results.tex both use \label{sec:results}; they
% are alternatives (a comprehensive vs a compact take) -- input only one.
% =====================================================================

\section{Results and Comparisons}
\label{sec:results}

We now report the measured cost of the two microcode kernels and compare each
against the best code that exists for it on the same silicon. The comparison is
deliberately unforgiving: for Keccak we bench against every scalar and SIMD
permutation from SUPERCOP that runs on this core, the eXtended Keccak Code
Package, and the OpenSSL assembly; for Curve25519 we bench against a formally
verified generator, a cryptographic superoptimiser, portable and
\texttt{\_\_uint128\_t} C, and the hand-written assembly of the primitive's own
designers. Every contender is measured on the same processor, under the same
harness, and---following the SUPERCOP autotuner---swept over the same grid of
compilers and optimisation levels with the best configuration retained for each.
The section is organised as Zhang et al.\ [TCHES 2025(1)] organise theirs: a
common setup and measurement preamble (\Cref{sec:results-setup}), then one
subsection per primitive (\Cref{sec:results-keccak,sec:results-curve}), each
leading with its comparison table and reading the numbers row by row, and closing
with a summary and future work (\Cref{sec:results-future}). We state the limits of
the advantage as plainly as the advantage itself.

% =====================================================================
\subsection{Experimental Setup and Measurement Methodology}
\label{sec:results-setup}
% =====================================================================

\paragraph{Platform.}
All measurements are obtained on an Intel Celeron N3350 (Goldmont
microarchitecture, family/model \texttt{0x506C9}), a low-power x86-64 core that is
\emph{out-of-order} and three-wide, with three integer ALU ports, running from a
single core pinned with \texttt{taskset\,-c\,0}. Its instruction set reaches
SSE4.2, AES-NI, and the SHA extensions, but---importantly for what follows---it
provides neither the BMI2 fused multiply \texttt{MULX} nor the ADX dual
carry-chain instructions \texttt{ADCX}/\texttt{ADOX} nor any of AVX/AVX2, so
native code on this part is confined to the single architectural carry flag and to
scalar or SSE-width vectors. Unless stated otherwise the governor is set to
\texttt{userspace} and the core frequency is pinned to its base of
$\sim\!1.09$\,GHz (nominal $1.10$\,GHz); the reason this matters is explained
below. The microcode patches are installed into the Goldmont patch RAM through the
red-unlock mechanism of \Cref{sec:ucode-layer}.

\paragraph{The cycle-counter pitfall, and why it dictates how we report.}
Timing on this part is subject to a pitfall that invalidates naive comparisons,
and we state it first because it shapes every number that follows. The
\texttt{rdtsc} instruction counts at a fixed reference rate of approximately
$1.1$\,GHz, whereas the core bursts to approximately $2.4$\,GHz under load, so the
\emph{same} work is reported as roughly $2.2\times$ more or fewer ticks depending
on the power state in force during the measurement. A cycle count captured in one
run and compared against a baseline captured in a separate run is therefore
meaningless. We take two precautions. First, we time the microcode and \emph{all}
of its competitors back to back within a single process, after a warm-up loop has
driven the core to a steady frequency, so that every contender observes the same
clock and the ratio between them is frequency-invariant. Second, for the absolute
cycle counts we additionally pin the core to its base frequency, so that a
reported tick approximates a true core cycle; a guard in the harness aborts the
run if any core is observed outside $\pm 6\%$ of base. The ratio is the robust
quantity throughout: the identical Keccak harness measured at burst frequency
reports approximately 874 against 942 cycles, the same ratio as the base-pinned
1912 against 2043, scaled by the $2.4/1.1$ power-state factor.

\paragraph{The compiler-and-optimisation sweep.}
A single compiler and flag choice is not a fair baseline, because the native
contenders are exactly the kind of code whose performance the compiler reshapes.
We therefore adopt the discipline of the SUPERCOP autotuner: each contender is
built and run under all 24 combinations of six compilers---\texttt{gcc}
11, 12, 13 and \texttt{clang} 14, 17, 18---and four optimisation
levels---\texttt{-O3}, \texttt{-O2}, \texttt{-Os}, \texttt{-O}---and the best
result for each contender is retained, with the winning configuration recorded.
The shared flags are \texttt{-fwrapv -fPIC -fPIE -march=native}, with
\texttt{-mtune=native} added for \texttt{gcc}. The microcode itself is a fixed
patch-RAM image and is compiler-independent, so its number is essentially flat
across the grid and serves as a sanity baseline; what the sweep genuinely tunes
are the C contenders, while the hand-written assembly baselines assemble
identically at every configuration.

\paragraph{Statistics and verification.}
For Curve25519 each contender is timed over 100 repetitions of a complete scalar
multiplication and we report the median. For Keccak, whose per-call cost is
smaller, each contender is timed in batches of 1000 permutations over 200
repetitions, all contenders interleaved within the single timing process; we
report both the median and a robust minimum that discards any batch reading below
half the median, since such readings can only arise from a power-state transition
straddling the measurement bracket. Correctness is a gate, not an afterthought:
the two Curve25519 field operations are checked against the fiat-crypto reference
over a large random corpus and a thousand-iteration self-chaining test, the
assembled scalar multiplication is checked against the RFC~7748 known-answer
vectors, and the Keccak permutation is checked against the Keccak team's published
$f(0)$ and $f(f(0))$ vectors, reproducing the canonical anchor
$\texttt{lane0}=\texttt{0xF1258F7940E1DDE7}$. No timing is reported for a kernel
that has not first passed its correctness gate.

% =====================================================================
\subsection{Keccak-f[1600]}
\label{sec:results-keccak}
% =====================================================================

\Cref{tab:results-keccak} reports the per-permutation cost of the looped
microcode Keccak against every permutation from the SUPERCOP
\texttt{crypto\_hash/keccakc1024} family that executes on this 64-bit Intel core,
together with the XKCP single-permutation kernel and the OpenSSL x86-64 assembly.
The variants that cannot run here are excluded and named for completeness: the
AMD-XOP kernel (\texttt{xopu24}) raises an illegal-instruction fault on a
GenuineIntel part, the other-ISA and 32-bit bit-interleaved kernels do not apply,
and the AVX2/AVX-512 OpenSSL paths require vector extensions Goldmont lacks. The
comparison is thus against the genuinely fastest code that runs on the hardware,
not a strawman.

\begin{table}[t]
\centering
\caption{Keccak-$f[1600]$ on the Goldmont N3350, best configuration per contender
over the 24-configuration compiler-and-optimisation sweep, measured base-pinned
and same-process. Cycles are per permutation; \texttt{this/$\mu$code} is the
contender's minimum divided by the microcode minimum, so a value above~1 means
the microcode is faster. A smaller cycle count is better. The microcode row is
the looped single-firing kernel of \Cref{sec:keccak-design}.}
\label{tab:results-keccak}
\small
\begin{tabular}{@{}llrrrl@{}}
\toprule
Contender & Method & Min & Median & \texttt{this/$\mu$code} & Best config \\
\midrule
\textbf{microcode} & looped, one firing, full-state resident & \textbf{1912} & 1917 & \textbf{1.000$\times$} & gcc-12 -O2 \\
\midrule
opt64lcu24        & C64: lane-compl., full unroll (24)       & 2043 & 2051 & 1.069$\times$ & gcc-11 -Os \\
openssl           & asm: x86-64 scalar (CRYPTOGAMS)          & 2047 & 2053 & 1.071$\times$ & gcc-11 -O  \\
x86\_64\_asm      & asm: hand-written, \texttt{ROL} (Van Keer) & 2064 & 2073 & 1.079$\times$ & gcc-11 -Os \\
xkcp\_g64lc       & XKCP plain-64: unroll, lane-compl.       & 2151 & 2160 & 1.125$\times$ & gcc-12 -O2 \\
opt64lcu6         & C64: lane-compl., unroll 6               & 2216 & 2218 & 1.159$\times$ & clang-17 -O3 \\
xkcp\_g64         & XKCP plain-64: unroll, no lane-compl.    & 2268 & 2270 & 1.186$\times$ & clang-18 -Os \\
simple            & reference C                              & 2370 & 2371 & 1.240$\times$ & gcc-11 -O3 \\
opt64u6           & C64: plain, unroll 6                     & 2394 & 2396 & 1.252$\times$ & gcc-12 -Os \\
sseu2             & SIMD: SSE2/SSSE3                         & 3047 & 3053 & 1.594$\times$ & clang-18 -Os \\
mmxu1             & SIMD: MMX                                & 3941 & 3953 & 2.061$\times$ & gcc-12 -O3 \\
opt64lcu24shld    & C64: \texttt{SHLD} rotation, unroll 24   & 10058 & 10069 & 5.260$\times$ & clang-17 -Os \\
x86\_64\_shld     & asm: hand-written, \texttt{SHLD}         & 10100 & 10112 & 5.282$\times$ & clang-17 -O \\
\bottomrule
\end{tabular}
\end{table}

\paragraph{The permutation as a whole.}
The looped microcode permutation completes in 1912 cycles, against 2043 for the
fastest scalar contender \texttt{opt64lcu24} and 2064 for the hand-written
\texttt{x86\_64\_asm}, a ratio of $0.936$ and an improvement of approximately
$6.8\%$. It is the fastest of every contender we measured. The three fastest
baselines---the lane-complemented, fully unrolled C kernel; the OpenSSL scalar
assembly; and the Van Keer rotate-based hand assembly---cluster within one percent
of one another at just above 2000 cycles, which is itself evidence that the
2043-cycle figure is a genuine ceiling for native code on this core rather than an
artefact of one implementation. The microcode sits clearly below that cluster.

\paragraph{The instruction the core hates.}
Two contenders stand a factor of five slower than the rest: both variants that
realise the $\rho$ rotations with the \texttt{SHLD} double-precision shift rather
than a rotate. On this Atom-class core \texttt{SHLD} is a heavily microcoded,
multi-cycle operation, and a kernel that issues it once per lane per round is
crippled by it. We include these rows not to flatter the microcode but because
they are instructive: the same source, differing only in how a rotation is
spelled, spans an order of magnitude, which is a concrete reminder that on this
microarchitecture the choice of primitive instruction, not the high-level
algorithm, dominates. The SIMD kernels (\texttt{sseu2}, \texttt{mmxu1}) are also
uncompetitive here, because Goldmont's narrow SSE-width vector units, with no AVX2,
handle the bitwise logic of Keccak poorly---an observation Zhang et al.\
[TCHES 2025(1)] make about their own RISC-V vector unit as well.

\paragraph{Compiler stability.}
Because a Keccak permutation runs entirely inside one microcode firing, the host
compiler touches only the thin native prologue that triggers the patch, so the
microcode is nearly invariant across the sweep. \Cref{tab:results-keccak-stab}
contrasts the competitive 64-bit contenders at their best and worst
configurations. The microcode varies by a single cycle across all 24
compiler-and-optimisation combinations, and the hand-written assembly---whose
generated code is likewise fixed---by four; but the fastest \emph{compiled}
baseline, \texttt{opt64lcu24}, swings by 213 cycles and attains its 2043-cycle
minimum only when one particular compiler happens to schedule it well, being some
$10\%$ slower at its worst. The microcode is thus not merely the fastest entry but
also the one whose performance a deployer could most reliably reproduce, since it
does not depend on the compiler finding a good schedule.

\begin{table}[t]
\centering
\caption{Stability across the 24-configuration sweep for the competitive 64-bit
contenders: the per-permutation cost at the best against the worst configuration,
and their spread (both read off the per-configuration minimum matrix). The
microcode and the two hand-assembly kernels are essentially compiler-invariant;
the fastest \emph{compiled} baseline \texttt{opt64lcu24} swings by over 200
cycles.}
\label{tab:results-keccak-stab}
\small
\begin{tabular}{@{}lrrr@{}}
\toprule
Contender & Best & Worst & Spread \\
\midrule
microcode (looped) & 1912 & 1913 & 1   \\
opt64lcu24         & 2043 & 2256 & 213 \\
openssl            & 2047 & 2050 & 3   \\
x86\_64\_asm       & 2064 & 2068 & 4   \\
xkcp\_g64lc        & 2151 & 2156 & 5   \\
opt64lcu6          & 2216 & 2222 & 6   \\
xkcp\_g64          & 2268 & 2273 & 5   \\
simple             & 2370 & 2378 & 8   \\
opt64u6            & 2394 & 2395 & 1   \\
\bottomrule
\end{tabular}
\end{table}

\paragraph{Where the advantage comes from, and its bound.}
The margin is modest by design, and the latency model of \Cref{sec:findings}
explains why. The kernel compiles to 108 static triads---a 26-triad prologue, a
56-triad looped body, and a 26-triad epilogue---which lies comfortably within the
128-triad patch RAM; because the body is looped rather than unrolled, the dynamic
count executed per permutation is $26 + 24\times56 + 26 = 1396$ triads. A single
isolated firing sustains $0.49$ cycles per triad, but that is a throughput figure
obtained by overlapping independent firings, whereas the looped permutation is a
strict dependency chain---round $N$ needs round $N{-}1$, with a backward branch on
each and no overlap---so it pays the true latency of approximately $1.37$ cycles
per triad, and $1396$ triads accordingly cost approximately 1912 cycles. The
microcode is thus latency-bound at roughly one triad per cycle, while the
hand-tuned scalar assembly, running on a three-wide out-of-order core, extracts
approximately three operations per cycle from the rotate and \texttt{ANDN}
instruction-level parallelism of Goldmont. We cannot close this gap with more
parallelism, because the chain $\theta\!\to\!\rho\!\to\!\pi\!\to\!\chi$ is
inherently sequential and the 128-triad cap precludes the unrolling that would
expose independent work across rounds. The entire advantage is therefore the
per-round work that native code \emph{cannot} perform: full register residency, so
that state input and output are paid once per permutation rather than once per
round; the borrowed thirty-second register that keeps the $\theta$-mixing lanes
$D$ resident; and the move-free $\chi$ stage that uses the freed registers as
scratch. \Cref{tab:results-keccak-opt} assembles the improvement from these
structural contributions, each measured back to back in the same harness. This is
the analogue of the ``normalized'' and CPI columns Zhang et al.\ report: it
exposes where the cycles go rather than only how many there are.

\begin{table}[t]
\centering
\caption{The Keccak improvement assembled from successive structural
optimisations, each expressed as the ratio to the naive residency-only baseline
(smaller is faster). Every step attacks the round's dependency chain; the
move-free $\chi$ is the single largest contributor.}
\label{tab:results-keccak-opt}
\small
\begin{tabular}{@{}clr@{}}
\toprule
Step & Change & Ratio \\
\midrule
0 & per-lane $D$ re-read from buffer, naive $\chi$ with auxiliary moves & $1.00\times$ \\
1 & $D$ resident in registers via the borrowed \texttt{RSP}, no base rebuild & $0.99\times$ \\
2 & move-free $\chi$ using the now-dead $D$-registers as scratch & $0.95\times$ \\
3 & depth-three $\theta$ parity tree and byte-indexed round constant & $0.93\times$ \\
\bottomrule
\end{tabular}
\end{table}

\paragraph{Cross-architecture context.}
Our claim is a same-silicon one: microcode against the best native code on the
\emph{same} core. For orientation only---and, as Zhang et al.\ themselves caution
about such comparisons, ``clearly unfair and provided for reference only''---we
note where optimised scalar Keccak-$f[1600]$ lands on other recent parts. Zhang
et al.\ [TCHES 2025(1)] report 1770 cycles on the dual-issue in-order XuanTie C908
with the RISC-V bit-manipulation extension, and cite Becker and Kannwischer's 1418
cycles on an ARM Cortex-A55 scalar pipeline, both exploiting a
bit-manipulation-rich instruction set and lazy-rotation techniques that Goldmont's
x86 ISA does not offer. These numbers are not comparable to our 1912 cycles in any
rigorous sense---different ISA, different core, different rotation and and-not
primitives---and we draw no performance conclusion from them. We include them
because they make the structural point precise: none of these native
implementations achieves full-state register residency, which is exactly the lever
the microcode uses and exactly what a register-limited ISA denies its compiler and
its assembly programmer.

% =====================================================================
\subsection{Curve25519 / X25519}
\label{sec:results-curve}
% =====================================================================

For Curve25519 the object placed in the patch RAM is not the scalar
multiplication but the two field operations beneath it, the multiplication and the
squaring modulo $p=2^{255}-19$ in the unsaturated radix-$2^{51}$ representation
(\Cref{sec:curve25519}). We therefore report at two levels: the field operation in
isolation, which is where the paper's claim is staked, and the complete scalar
multiplication, which orients that claim within the state of the art.

\paragraph{The field operation in isolation.}
Timed on their own in a dependent chain---the way the ladder issues them---the
microcode field operations are faster than the fiat-crypto reference against which
the superoptimiser is itself measured. The squaring completes in roughly
37 cycles, about $26\%$ faster than fiat, and the multiplication in roughly
58 cycles, about $3\%$ faster. The asymmetry is intrinsic to the kernels: squaring
issues fifteen partial products against multiplication's twenty-five, so it has
proportionally more slack into which the concurrent carry chains of
\Cref{sec:curve25519} can be packed, whereas the multiplication is nearer the
multiplier-throughput bound where the layer's carry-concurrency advantage has less
room to act. Because both the Montgomery ladder and the Fermat inversion are
squaring-dominated, it is the larger squaring advantage that propagates into the
end-to-end figure.

\paragraph{The cleanest comparison: the same framework, only the field ops
change.} The most confounding-free way to attribute an end-to-end difference to
the field operations is to hold everything else constant. We do this by carrying
the microcode field operations into the Bernstein--Schwabe \texttt{amd64-51}
framework, leaving its ladder, inversion, conditional swap, and (un)packing
exactly as its authors wrote them, and swapping only \texttt{fe25519\_mul} and
\texttt{fe25519\_square} for the microcode. Across the 24-configuration sweep the
microcode field operations make the same framework a geometric mean of
$1.084\times$ faster than the authors' own hand-written assembly field
operations---about $8\%$---so the advantage is not an artefact of our own ladder
and survives transplantation into an independently engineered scalar
multiplication. That such a chained accumulation is where the layer helps at all
is itself a consequence of the platform: lacking the ADX extension, native
Goldmont code cannot overlap two carry chains through \texttt{ADCX}/\texttt{ADOX}
and must serialise on the single architectural carry flag, whereas the microcode
advances three accumulation chains at once, each on its own private flag domain
(\Cref{sec:findings}).

\paragraph{End to end.}
\Cref{tab:results-curve-e2e} reports the best per-contender result across the
sweep, in median cycles per complete scalar multiplication, alongside the
geometric mean over all 24 configurations of each contender's ratio to our inlined
microcode ladder. The inlined microcode ladder completes in about 300\,000 cycles
and is the fastest of every $5\times51$ contender we measured: it is a geometric
mean of $1.25\times$ faster than fiat-crypto, $1.30\times$ faster than
hand-written \texttt{\_\_uint128\_t} C, $1.28\times$ faster than the CryptOpt
superoptimiser, $1.31\times$ faster than \texttt{donna\_c64}, and $1.17\times$
faster than the hand-written \texttt{amd64-51} assembly stack. Against a formally
verified generator, a cryptographic superoptimiser, expert portable C, and expert
hand assembly, all at or above the ISA, the below-ISA microcode is the fastest
$5\times51$ X25519 on this core.

\begin{table}[t]
\centering
\caption{End-to-end X25519 on the Goldmont N3350, best configuration per contender
over the 24-configuration sweep, in median cycles. The last column is the
geometric mean over all 24 configurations of (contender) $\div$
(ours/microcode-inline); above~1 means our microcode is faster, below~1 means it
is slower. The inlined microcode ladder is the fastest $5\times51$ implementation
but is overtaken by the hand-written $4\times64$ saturated assembly, and the
microcode port \emph{into} that $4\times64$ representation is the slowest entry by
a wide margin.}
\label{tab:results-curve-e2e}
\small
\begin{tabular}{@{}llrr@{}}
\toprule
Contender & Representation and backend & Median (cyc) & geomean vs ours \\
\midrule
amd64-64 / asm       & $4\times64$ saturated, hand asm             & 271\,326 & $0.891\times$ \\
\textbf{ours / ucode-inline} & $5\times51$, microcode (inlined ladder) & \textbf{299\,967} & $1.000\times$ \\
amd64-51 / ucode     & $5\times51$, microcode (B--S framework)     & 314\,416 & $1.080\times$ \\
donna\_c64           & $5\times51$, portable C                     & 335\,744 & $1.307\times$ \\
ours / fiat          & $5\times51$, fiat-crypto C                  & 354\,851 & $1.253\times$ \\
amd64-51 / asm       & $5\times51$, hand asm                       & 356\,038 & $1.170\times$ \\
ours / hand-C        & $5\times51$, hand-written \texttt{\_\_uint128\_t} C & 379\,433 & $1.300\times$ \\
ours / cryptopt      & $5\times51$, CryptOpt asm                   & 380\,361 & $1.277\times$ \\
amd64-64 / ucode     & $4\times64$ saturated, microcode            & 511\,017 & $1.773\times$ \\
\bottomrule
\end{tabular}
\end{table}

\paragraph{Where the advantage ends.}
The fastest scalar multiplication on this core is nonetheless not ours. It is the
Bernstein--Schwabe \texttt{amd64-64} implementation, which uses the saturated
$4\times64$ representation in radix $2^{64}$ and completes in 271\,326 cycles, a
geometric mean of $1.12\times$ ahead of the inlined microcode ladder---about
$11\%$. Its advantage has nothing to do with carry handling, at which the
microcode excels; it is one of algorithm and of native execution. The saturated
$4\times64$ multiplication performs sixteen partial products against our
twenty-five, on a core whose multiplier, not its carry logic, is the scarce
resource; and being ordinary native code it pays none of the per-firing dispatch
tax that our hooked operations incur on the order of two and a half thousand times
per scalar multiplication. The honest headline is therefore precise: the microcode
is the fastest \emph{$5\times51$} X25519 we measured and the fastest field backend
under a fixed representation, but it is not the fastest X25519, because it cannot
adopt the better algorithm.

\paragraph{The representation that suits the layer worst.}
The last row of \Cref{tab:results-curve-e2e} makes the boundary sharp. When we
encode the field operations themselves in the $4\times64$ representation, the
microcode port runs in 511\,017 cycles---a geometric mean of about $1.77\times$
slower than our $5\times51$ microcode and nearly twice the cost of the native
$4\times64$ assembly. The cause is exactly the flag structure of
\Cref{sec:findings}. A saturated $4\times64$ multiplication produces long,
full-width carry chains, and the add-with-carry instruction those chains want
reads the architectural flag domain that is frozen on entry to a patch and cannot
be threaded through a sequence of dependent additions inside the microcode. The
four-deep saturated carry must therefore be emulated by serial \texttt{SETCC}
operations, reintroducing precisely the serialisation the $5\times51$ kernel's
three parallel carry chains were built to avoid. The representation that suits the
native core best is thus the one that suits the layer worst, and the
representation at which the layer excels does not yield the fastest algorithm
overall. This is the crisp form of the headroom result: descending to microcode
buys a single-digit-percent constant factor on a \emph{fixed} algorithm, not the
freedom to run a better one.

% =====================================================================
\subsection{Summary and Future Work}
\label{sec:results-future}
% =====================================================================

\paragraph{Summary.}
Both results have the same shape. A complete, known-answer-verified 24-round
Keccak-$f[1600]$ runs entirely in Goldmont microcode, looped within a single
firing in 108 triads, and is approximately $7\%$ faster than the fastest
hand-tuned scalar assembly on this microarchitecture---the fastest of every
runnable contender, and the most stable across compilers. A hand-scheduled
microcode encoding of the Curve25519 $5\times51$ field operations is the fastest
such field backend we measured, ahead of an optimising compiler, a formally
verified generator, a cryptographic superoptimiser, and expert hand assembly, and
the advantage survives transplantation into the Bernstein--Schwabe framework,
where it is a geometric mean of about $8\%$ faster than the authors' own assembly.
In both cases the advantage is structural---concurrent private-flag carry chains
fused into single triads for the field arithmetic, whole-state register residency
with move-free $\chi$ for the permutation---and in both cases it is bounded, by the
128-triad patch RAM, the per-firing dispatch tax, and one-triad-per-cycle issue,
not by anything the compiler leaves on the table. Where the layer loses it loses
for a reason we can name: the faster X25519 belongs to a $4\times64$ saturated
algorithm that the frozen single-flag model makes the layer worst at, and that the
patch RAM cannot in any case hold.

\paragraph{Future work.}
Several levers remain. For X25519, replacing the Fermat inversion with a
Bernstein--Yang constant-time inversion is projected to recover a further
$1$--$2.5\%$ end-to-end and is gated on a short constant-time probe. The
same-process interleaved harness used for Keccak should be extended to the
X25519 field-operation comparison so that its claim rests on a same-process ratio
rather than cross-run minima, and the reported statistics should be broadened from
median and minimum to include an inter-quartile range. More broadly, the method
generalises: we have applied the same generator-and-simulator methodology to
several further primitives (P-256, secp256k1, poly1305), which a fuller treatment
could report; and the Keccak kernel is directly reusable as the permutation inside
hash-based signatures. Finally, the security surface this work opens---that a
primitive which can be \emph{implemented} in microcode can equally be
\emph{subverted} there, invisibly to any ISA-level audit---deserves the separate
treatment we give it in \Cref{sec:ucode-layer} and beyond.
